\section{Exact section growth from the classical abelian-variety index theorem}
\label{sec:SGV}

\status{Seven detailed claims supply the complete original
\texttt{secM\_finite\_and\_growth} output.  This is a classical
replacement proof using PBC26--PBC27 and the explicitly stated Milne
inputs below.  It installs the independently reviewed comparison
argument without changing any earlier proof.}

The point of this chapter is to prove the precise numerical input
needed by GSL19 on the original section modules.  We choose one very
ample power of the given line bundle and compute all its positive
section dimensions exactly.  In particular, there is no need for this
consumer to reconstruct the more general theorem about growth of
filtered polynomial images.  That general theorem remains useful in
other settings; it is not claimed here.

\subsection*{The original objects and the explicit classical input}

Fix an algebraically closed field $K$, with no restriction on its
characteristic, and put $S=\operatorname{Spec}K$.  Let $P$ be an
original $\operatorname{PolAV}(S)$, and write
$A=P.A$, $L=P.L$, $g=P.g$ and $a:A\to S$ for the structural map.
The record gives a group scheme whose structural map is smooth, proper
and geometrically connected; it specifies smooth relative dimension
$g$.  The datum $L$ is an actual line bundle with the original
rigidification, symmetry, relative ampleness and $\operatorname{KTriv}$
fields.  We keep this entire record.  The proof below needs its
ampleness, not an additional identification of its translation kernel
with a Picard-scheme kernel.

Let $T=\operatorname{baseOver}(S)$ be the identity object over $S$.
The actual product in the definition of the sections is
\[
 X=(T\otimes A).\mathrm{left},\qquad
 \pi=(\operatorname{snd}(T,A)).\mathrm{left}:X\longrightarrow A,
 \qquad
 f=(\operatorname{fst}(T,A)).\mathrm{left}:X\longrightarrow S.
\]
Put $R_0=\Gamma(T,\mathcal O_T)$ and, for every integer $j$, put
\[
 V_j=\operatorname{SecM}(A,L,j,T).
                                                        \tag{SGV.1}
\]
These are exactly the source's modules, including its scalar action
through the first projection.

We use three stated parts of classical abelian-variety theory from
J.~S.~Milne, \emph{Abelian Varieties}, corrected 2022 version of the
1986 Storrs article,
\url{https://www.jmilne.org/math/xnotes/AVs.pdf}.
For an abelian variety $B$ over a field and an invertible sheaf $M$,
write $K_M$ for its translation stabilizer, with Milne's reduced
closed-subscheme convention.  Proposition~9.1 on printed page~15
implies that an ample $M$ with a nonzero section has
$\dim K_M=0$.  Theorem~13.3(c) on printed page~25 says that this last
condition concentrates the nonzero cohomology in one degree.
Theorem~13.3(b) on the same page gives
\[
 \chi(B,\mathcal O_B(D))=\frac{(D^{\dim B})}{(\dim B)!}
                                                        \tag{SGV.2}
\]
for a Cartier divisor $D$.  These results hold in arbitrary
characteristic.  They are retained as published classical inputs;
their proofs are not being attributed to this chapter.  We do not
use Theorem~13.3(a), the degree formula for a polarization.

The elementary scheme dictionary uses smoothness over a field implying
regularity, regular local rings being domains, and the fiber-dimension
meaning of smooth relative dimension.  References are the Stacks
Project tags \nolinkurl{04QM}, \nolinkurl{00NN} and
\nolinkurl{01V4}; the component and connectedness argument is given
explicitly below.  The two inherited human proofs used are PBC26
(very ample powers on every actual affine test) and PBC27 (finiteness
of the actual sections, including every integer power).

% BEGIN SGV1
\paragraph{SGV1: the actual identity product is an abelian variety of dimension \(g\).}
\label{claim:SGV1}
The specified projection $\pi:X\to A$ is an isomorphism over $S$.
With the transported group law, $X$ is a geometrically integral
abelian variety, and its dimension is the given integer $g=P.g$.

\begin{proof}
The defining square is Cartesian with lower map
$\operatorname{id}_S$.  Its universal property supplies
$\beta:A\to X$ from the pair $(a,\operatorname{id}_A)$.
The second projection gives $\pi\beta=\operatorname{id}_A$.
Both projections of $\beta\pi$ agree with the corresponding
projections of $\operatorname{id}_X$, so the uniqueness part of the
same universal property gives $\beta\pi=\operatorname{id}_X$.
Thus these are mutually inverse maps over the specified base.  In
particular, transporting multiplication, identity and inverse through
$\pi$ gives the existing group scheme in this identity-product model.

The original $\operatorname{isAb}$ fields imply that $X\to S$ is
smooth, proper and geometrically connected.  Properness gives finite
type and separatedness over the field, so $X$ is Noetherian.  Its local
rings are regular, hence domains.  Two distinct irreducible components
cannot meet: at an intersection point they would give distinct
minimal primes of the same local ring.  More explicitly, take an
affine neighborhood of that point.  The distinct component primes
contained in its point prime remain distinct on localization; if an
element of one entered the other after localization, multiplying it
by an element outside the point prime would contradict primality.
A domain has only one minimal prime, giving the contradiction.

There are finitely many irreducible components.  They are disjoint
closed sets, hence each is also open.  Connectedness and the identity
point therefore imply that there is exactly one component.  Every
local ring is reduced, so $X$ is reduced as well and is integral.

The argument applies after every extension of $K$: smoothness,
properness, geometric connectedness and the identity section all
persist.  Thus $X$ is geometrically integral.  It is a complete
connected group variety, which is the abelian-variety convention in
Milne's article.  Finally the original smooth relative dimension $g$
transports to $f$.  Over a field every component of its unique fiber
has dimension $g$; this fiber is $X$ itself.  Hence the dimension in
(SGV.2) is exactly $P.g$.
\end{proof}
% END SGV1

% BEGIN SGV2
\paragraph{SGV2: the literal section module and its scalar field.}
\label{claim:SGV2}
Let $\mathcal L$ be the invertible sheaf defined by the original
chart datum $L$.  For every $j\in\mathbb Z$ there is a natural
coordinate-preserving identification
\[
 V_j\simeq H^0\bigl(X,(\pi^*\mathcal L)^{\otimes j}\bigr).
                                                        \tag{SGV.3}
\]
It is linear after transporting scalars through the canonical field
isomorphism $e:K\simeq R_0$.  Consequently the corresponding finite
dimensions over $K$ and $R_0$ agree.

\begin{proof}
The original datum consists of opens $U_i$ and transition units
$g_{ij}$ satisfying the cover and cocycle identities.  Glue
$\mathcal O_{U_i}$ by these units.  The resulting sheaf is invertible,
and a global section is exactly a tuple $a_i$ with
$a_i=g_{ij}a_j$ on each overlap, with the restrictions understood.
The gluing construction preserves sums and multiplication by
functions on the base.

By definition $\operatorname{LPow}(A,L,j,T)$ is the
singleton-factor $\operatorname{LinComb}$ of $L$ along $\pi$,
with exponent $j$.  Its charts are indexed by functions
$\operatorname{Fin}(1)\to L.\iota$.  Evaluation at $0$ identifies
these with the original chart indices.  The chart opens are
$\pi^{-1}U_i$, and the transition units are $(\pi^*g_{ij})^j$.
They are exactly those of $(\pi^*\mathcal L)^{\otimes j}$,
using the dual bundle when $j<0$.

The submodule $\operatorname{SecM}$ imposes precisely these overlap
equations.  The source's maps $\operatorname{SecM.toSec}$ and
$\operatorname{SecM.ofSec}$ leave every chart coordinate unchanged,
so the preceding description proves (SGV.3) as a bijection of the
actual sections, not merely as an equality of dimensions.

For $b\in R_0$ the source action on a coordinate over an open
$W\subseteq X$ is multiplication by $f^*b|_W$.
The canonical isomorphism
$e:K\simeq\Gamma(\operatorname{Spec}K,\mathcal O)$ identifies
this action with the ordinary $K$-action on the right of (SGV.3).
Thus the identification preserves addition and the specified scalar
action.  A basis over one field transports to a basis over the other,
which proves the dimension assertion.  No sign restriction on $j$
was used.
\end{proof}
% END SGV2

% BEGIN SGV3
\paragraph{SGV3: one fixed very ample power, with nonzero sections in every positive power.}
\label{claim:SGV3}
There is an integer $N\ge1$ such that
$H=(\pi^*\mathcal L)^{\otimes N}$ is very ample.
The integer $N$ is chosen before any positive $t$.  For every
$t\ge1$, the sheaf $H^{\otimes t}$ is ample and has a nonzero
global section.  Its actual coordinate sections identify with $V_{Nt}$.

\begin{proof}
Apply PBC26 to the displayed identity Cartesian square, with $R=K$.
It supplies $N\ge1$ and actual sections $s_0,\ldots,s_m$ of
$\operatorname{pw}(L.\operatorname{pull}(\pi),N)$ satisfying
$\operatorname{VeryAmpleBy}(f,\operatorname{pw}(\pi^*L,N),s)$.
The associated invertible sheaf is $H$.

Here the original coordinate form of very ampleness gives classical
very ampleness directly.  The opens $D_i$ on which $s_i$ is nonzero
cover $X$, and $s_i$ trivializes $H$ on $D_i$.  The supplied ratio
functions satisfy $r_j=s_j/s_i$ there, so $r_i=1$.
These ratios define the morphism $X\to\mathbb P^m_K$ given by the
sections.  Its inverse image of the $i$th standard affine chart is
exactly $D_i$.  The specified polynomial evaluation map to
$\Gamma(D_i,\mathcal O)$ is surjective.  Since it sends its $i$th
variable to $1$, it factors through that relation and gives a
surjective map from the polynomial ring in the remaining $m$ chart
coordinates.  Each map to the standard affine chart is therefore a
closed immersion.  Closed immersion is local on the target, so they
give a closed immersion into $\mathbb P^m_K$ whose pullback of
$\mathcal O(1)$ is $H$.

Choose an $i$ for which $s_i$ is nonzero at the identity point of $X$.
Such a choice is possible because the $D_i$ cover.  For every $t\ge1$,
$s_i^{\otimes t}$ is nonzero at that same point: it is the tensor
power of a nonzero vector in a one-dimensional $K$-space.
This proves $H^0(X,H^{\otimes t})\ne0$.  Composition of the displayed
closed immersion with the degree-$t$ Veronese embedding proves that
$H^{\otimes t}$ is very ample and hence ample, in any characteristic.

Lastly, iterating the power-chart operation gives transition units
$((\pi^*g_{ij})^N)^t=(\pi^*g_{ij})^{Nt}$.  Identifying the singleton
chart indices as in SGV2 gives the actual coordinate correspondence
with $V_{Nt}$.  The restriction maps and scalar action are unchanged.
This is the coordinate content of the source's mutually inverse
$\operatorname{toPw}$ and $\operatorname{ofPw}$ maps.
\end{proof}
% END SGV3

% BEGIN SGV4
\paragraph{SGV4: concentration in degree zero for every positive power of \(H\).}
\label{claim:SGV4}
For every integer $t\ge1$,
\[
 H^q(X,H^{\otimes t})=0\quad(q>0),\qquad
 h^0(X,H^{\otimes t})=\chi(X,H^{\otimes t}).
                                                        \tag{SGV.4}
\]

\begin{proof}
By SGV1 the scheme is an abelian variety over the original field.
By SGV3 the sheaf $H^{\otimes t}$ is ample and has a nonzero
global section.  Milne's Proposition~9.1 therefore gives dimension
zero for its translation stabilizer.  Theorem~13.3(c) says that its
nonzero cohomology occurs in exactly one degree.  Since degree zero
is already nonzero, this degree must be zero.  Thus every higher
cohomology group vanishes, and the alternating sum defining the Euler
characteristic reduces to $h^0$.

This uses only the dimension-zero condition on the reduced translation
stabilizer in Milne's statement.  It does not require identifying it
with the original $\operatorname{KTriv}$ functor, and it does not use a
separability or prime-to-characteristic assumption on $t$.
\end{proof}
% END SGV4

% BEGIN SGV5
\paragraph{SGV5: the exact positive-dimensional growth formula.}
\label{claim:SGV5}
If $g>0$, then for every integer $t\ge1$,
\[
 h^0(X,H^{\otimes t})=t^g h^0(X,H).
                                                        \tag{SGV.5}
\]

\begin{proof}
Use the nonzero section $s_i$ chosen in SGV3.  On the integral scheme
$X$ it defines an effective Cartier divisor $D$ with
$H\simeq\mathcal O_X(D)$.  To see the Cartier condition explicitly,
its coordinate on any nonempty trivializing open is not zero: a zero
coordinate would force the section to vanish on that open, then on
every other chart through nonempty overlaps and the transition units.
The local coordinate rings are domains, so these nonzero coordinates
are nonzerodivisors.  Their principal ideals fit together because the
coordinates differ by units.  This is exactly the asserted effective
Cartier divisor, with the original bundle as its associated sheaf.

Tensor powers now give $H^{\otimes t}\simeq\mathcal O_X(tD)$.
By SGV4, Milne's Riemann--Roch formula and multilinearity of the
intersection product,
\[
 \begin{aligned}
 h^0(X,H^{\otimes t})
   &=\chi(X,\mathcal O_X(tD))
     =\frac{((tD)^g)}{g!}\\
   &=t^g\frac{(D^g)}{g!}
     =t^g\chi(X,H)
     =t^g h^0(X,H).
 \end{aligned}
\]
All intersections and Euler characteristics here are integer-valued;
the displayed fractions are interpreted in $\mathbb Q$.  We never
divide by $g!$ inside $K$.  Therefore this calculation remains valid
when the characteristic divides $g!$ or $t$.
\end{proof}
% END SGV5

% BEGIN SGV6
\paragraph{SGV6: the zero-dimensional case.}
\label{claim:SGV6}
If $g=0$, then $X\simeq\operatorname{Spec}K$, and (SGV.5) holds
for every $t\ge1$, with $h^0(X,H)=1$.

\begin{proof}
An integral zero-dimensional scheme has just its generic point:
any further point would give a strict chain of irreducible closed
subsets of length one.  Thus an affine neighborhood of the unique
point of $X$ is all of $X$.  Its coordinate ring is a domain of
dimension zero, hence a field.  It is a finite-type $K$-algebra by
SGV1, so Zariski's lemma makes it a finite algebraic extension of $K$.
Algebraic closedness gives the equality with $K$.

An invertible sheaf on $\operatorname{Spec}K$ is a one-dimensional
$K$-vector space.  Every tensor power again has one-dimensional
sections.  Since $t^0=1$, this proves the formula, including $t=1$,
without using intersection theory in dimension zero.
\end{proof}
% END SGV6

% BEGIN SGV7
\paragraph{SGV7: the complete original section-finiteness and growth output.}
\label{claim:SGV7}
For every original $P$ over every algebraically closed field $K$,
$V_j$ is finite over $R_0$ for every $j\in\mathbb Z$.
There exist integers $N,d\ge1$, chosen before $t$, such that
\[
 \dim_{R_0}V_{Nt}=d\,t^g\qquad(t\ge1).
                                                        \tag{SGV.6}
\]
In particular, the full original conclusion is supplied:
\[
 \begin{gathered}
  (\forall j\in\mathbb N)\quad V_j\text{ is finite over }R_0,\\
  \exists N,c,C\in\mathbb N,\quad N,c\ge1,\quad
  \forall t\in\mathbb N,\ t\ge1:\quad
  t^g\le c\dim_{R_0}V_{Nt},\qquad
  \dim_{R_0}V_{Nt}\le C t^g.
 \end{gathered}
                                                        \tag{SGV.7}
\]

\begin{proof}
PBC27 gives finiteness of the actual $V_j$ for every integer $j$.
Choose the single $N$ from SGV3 and put $d=h^0(X,H)$.
This is a finite integer by SGV2 and the $j=N$ instance of PBC27.
It is positive because SGV3 gives a nonzero section.

SGV5 and SGV6 give (SGV.5) in every dimension.  The actual
power-coordinate comparison in SGV3 identifies the left side with
$V_{Nt}$, and SGV2 transports its dimension to the precise coefficient
field $R_0$.  This proves (SGV.6).  Since $d\ge1$, the choices
$c=1$ and $C=d$ prove both inequalities in (SGV.7).
Every choice precedes the universal quantifier over $t$.
The integers may depend on the given $K$ and $P$, as in the original
source; no uniformity across different polarized varieties is claimed.
There is no characteristic restriction on $K$, $N$ or $t$.
\end{proof}
% END SGV7

\subsection*{What this supplies next}

Equation (SGV.7) is exactly the finite-sections and growth output of
\nolinkurl{CubeAssembly.sections_growth}, and hence the growth input
called \texttt{hgrow} by
\nolinkurl{finrank_sections_lineBundlePow_eq_pow}.  The source obtains
it from several general geometric and polynomial-growth suppliers.
The present argument reaches the same original output directly using
the explicit classical input, while preserving the actual identity
product, section coordinates and scalar field.

For GSL19, SGV1 supplies the needed integrality of the original
geometric fiber and SGV7 supplies its fourth input.  The rank-one
statement $\dim V_1=1$ and the actual torsion-indexed section
equivalences remain to be proved.  In particular, the constant $d$ in
(SGV.6) is not declared to be $N^g$; establishing all ranks $j^g$
requires those further inputs.  The full all-cover eventual Serre
vanishing theorem remains supplied by SHV17, and its homogeneous-polynomial
companion remains supplied by SHV21.  The later descent theorem giving
all-cover vanishing for every positive power $L^n$ is a separate input,
which SGV does not supply.  Cube and square generators, translated-power
generators, the global $\operatorname{SecBaseChangeSt}$,
$\operatorname{SecBaseChangeAnySt}$, $\operatorname{SecTwistSt}$ and
moduli objects also remain outside this chapter.  No selected whole
obligation is closed here.

\subsection*{Exact original source correspondence}

All source paths below use commit
\nolinkurl{7cbf7a0c32fca1badf05931879a37c3f8a2b03b4} of
\url{https://github.com/CocoSgt/riemann-hypothesis-lean}.
Let $\mathcal A$ denote
\path{src/Development/AbelianVarietyTheory/Proofs/AbelianSchemeLawProofs/},
and $\mathcal R$ denote
\path{src/ReciprocityTheory/ReductiveGroups/ArithmeticGeometry/}.
\begin{enumerate}
 \item The actual record is
 $\mathcal R$\path{ExtensionDescent/NonzeroLocus.lean}, lines59--73;
 its abelian-scheme fields are in
 $\mathcal R$\path{SemiabelianSchemes/AbelianSchemeQuotients/IsAbelianScheme.lean},
 lines13--17.
 \item The line-bundle charts are
 $\mathcal R$\path{SemiabelianSchemes/CubicalStructures/CubeSymmetricSwapFirst.lean},
 lines17--78.  The literal power, action and section module are
 $\mathcal R$\path{ExtensionDescent/ThetaPt.lean}, lines53--57, and
 $\mathcal R$\path{ExtensionDescent/ThetaVV.lean}, lines18--60.
 \item The target and actual power-coordinate maps are in
 $\mathcal A$\path{FiniteOfNotLinearIndependent.lean}, lines280--350;
 its exported output is
 $\mathcal A$\path{CubeTheoremAndSectionsFiniteGrowthAbelianVariety.lean},
 lines38--43.
 \item The exact next rank consumer is
 $\mathcal A$\path{FinrankSectionsLineBundlePowEqPow.lean}, lines228--260.
 The separate inputs at the later global join appear in
 $\mathcal A$\path{LawsOfH0EqOneClosedFields.lean}, lines64--93.
\end{enumerate}
